\documentclass[11pt,a4paper]{article}
\usepackage[margin=1.8cm,headheight=48pt,headsep=12pt,footskip=22pt]{geometry}
\usepackage{fontspec}
\setmainfont{Noto Serif}
\setsansfont{Noto Sans}
\setmonofont{DejaVu Sans Mono}
\usepackage{amsmath,amssymb,mathtools}
\usepackage{enumitem,tabularx,booktabs,array,multicol}
\usepackage{xcolor,graphicx,tikz,fancyhdr,lastpage}
\usepackage{listings,tcolorbox,hyperref}
\hypersetup{colorlinks=true,linkcolor=blue!60!black,urlcolor=blue!60!black}
\definecolor{hcmusblue}{RGB}{28,55,111}
\definecolor{lightblue}{RGB}{237,242,250}
\definecolor{codebg}{RGB}{247,248,250}
\setlength{\parindent}{0pt}
\setlength{\parskip}{4pt}
\setlist[enumerate]{leftmargin=*,itemsep=5pt,topsep=4pt}
\setlist[itemize]{leftmargin=*,itemsep=2pt,topsep=2pt}
\lstset{language=C++,basicstyle=\ttfamily\footnotesize,backgroundcolor=\color{codebg},frame=single,rulecolor=\color{black!15},breaklines=true,showstringspaces=false,columns=fullflexible,keywordstyle=\color{hcmusblue}\bfseries,commentstyle=\color{green!40!black},numbers=left,numberstyle=\tiny\color{black!45},xleftmargin=1.5em,framexleftmargin=1.2em}
\pagestyle{fancy}
\fancyhf{}
\fancyhead[L]{\includegraphics[height=1.05cm,keepaspectratio]{../../04_Nhan_dien/hcmus_logo.png}\hspace{0.45em}\begin{minipage}[c]{0.64\textwidth}\textbf{TRƯỜNG ĐẠI HỌC KHOA HỌC TỰ NHIÊN}\\[-1pt]\small ĐHQG TP. HỒ CHÍ MINH\end{minipage}}
\fancyhead[R]{\small\textcolor{hcmusblue}{TÀI LIỆU LUYỆN TẬP}}
\fancyfoot[L]{\scriptsize Biên soạn độc lập - Không phải đề thi chính thức}
\fancyfoot[R]{\scriptsize Trang \thepage}
\renewcommand{\headrulewidth}{0.5pt}
\renewcommand{\footrulewidth}{0.3pt}
\newcommand{\examtitle}[4]{\begin{center}{\Large\bfseries\color{hcmusblue}#1}\par\vspace{3pt}{\large\bfseries #2}\par\vspace{3pt}\small #3\quad\textbullet\quad #4\end{center}\vspace{5pt}\hrule\vspace{8pt}}
\newcommand{\studentline}{\begin{tabularx}{\textwidth}{X X}\textbf{Họ và tên:} \hrulefill & \textbf{MSSV:} \hrulefill\\[7pt]\textbf{Lớp:} \hrulefill & \textbf{Ngày làm:} \hrulefill\end{tabularx}\vspace{8pt}}
\newcommand{\points}[1]{\hfill\textbf{(#1 điểm)}}
\newtcolorbox{notebox}{colback=lightblue,colframe=hcmusblue!55,boxrule=0.5pt,arc=1.5mm,left=2mm,right=2mm,top=1mm,bottom=1mm}
\begin{document}
\small
\examtitle{ĐÁP ÁN ĐỀ LUYỆN TẬP 02 - CƠ SỞ LẬP TRÌNH}{CƠ SỞ LẬP TRÌNH - HCMUS (theo lịch học tuần 1--5)}{Lời giải chi tiết và mã C++11}{Giữa kỳ 1 - HK1 2026--2027}

\section*{Đáp án và lời giải - Đề 02}
\begin{enumerate}
\item \begin{minipage}[t]{\linewidth}
\begin{minipage}[t]{\linewidth}
Điều kiện tam giác là $a+b>c$, $a+c>b$, $b+c>a$. Sau đó phân loại:
\begin{lstlisting}
if (a <= 0 || b <= 0 || c <= 0 ||
    a + b <= c || a + c <= b || b + c <= a) {
    cout << "Khong phai tam giac";
} else if (a == b && b == c) {
    cout << "Tam giac deu";
} else if (a == b || a == c || b == c) {
    cout << "Tam giac can";
} else {
    cout << "Tam giac thuong";
}
\end{lstlisting}
\end{minipage}
\end{minipage}
\item \begin{minipage}[t]{\linewidth}
\begin{minipage}[t]{\linewidth}
Mã C++11 tham khảo:
\begin{lstlisting}
int gcd(int a, int b) {
    while (b != 0) {
        int r = a % b;
        a = b;
        b = r;
    }
    return a;
}
\end{lstlisting}
Lặp cho đến khi số dư bằng 0; khi đó $a$ là ước chung lớn nhất. Ví dụ \texttt{gcd(18,12)}: $(18,12)\to(12,6)\to(6,0)$, kết quả 6.
\end{minipage}
\end{minipage}
\item \begin{minipage}[t]{\linewidth}
\begin{minipage}[t]{\linewidth}
Mã C++11 tham khảo:
\begin{lstlisting}
int maxValue(int a, int b) {
    return (a > b) ? a : b;
}

double maxValue(double a, double b) {
    return (a > b) ? a : b;
}
\end{lstlisting}
Hai lời gọi lần lượt chọn phiên bản nhận \texttt{int} và phiên bản nhận \texttt{double}; kết quả in ra là \texttt{5 2.5}.
\end{minipage}
\end{minipage}
\item \begin{minipage}[t]{\linewidth}
\begin{minipage}[t]{\linewidth}
Mã C++11 tham khảo:
\begin{lstlisting}
void reverseRange(int a[], int left, int right) {
    while (left < right) {
        int t = a[left];
        a[left] = a[right];
        a[right] = t;
        ++left;
        --right;
    }
}

void rotate(int a[], int n, int k) {
    k %= n;
    if (k < 0) k += n;
    if (k == 0) return;
    reverseRange(a, 0, n - 1);
    reverseRange(a, 0, k - 1);
    reverseRange(a, k, n - 1);
}
\end{lstlisting}
Ví dụ $[1,2,3,4,5]$ quay phải 2 vị trí thành $[4,5,1,2,3]$. Chuẩn hóa $k$ đưa số bước về đoạn $[0,n-1]$.
\end{minipage}
\end{minipage}
\item \begin{minipage}[t]{\linewidth}
\begin{minipage}[t]{\linewidth}
Mã C++11 tham khảo:
\begin{lstlisting}
void printEvenCountPerRow(const int a[][MAX], int R, int C) {
    for (int i = 0; i < R; ++i) {
        int count = 0;
        for (int j = 0; j < C; ++j)
            if (a[i][j] % 2 == 0) ++count;
        cout << count << '\n';
    }
}
\end{lstlisting}
Vòng ngoài duyệt từng hàng; vòng trong đếm phần tử chẵn của hàng đó.
\end{minipage}
\end{minipage}
\end{enumerate}

\end{document}
